\documentclass[11pt,a4paper]{article}
\usepackage[utf8]{inputenc}
\usepackage[T1]{fontenc}
\usepackage{times}
\usepackage{geometry}
\geometry{margin=2.5cm}
\usepackage{booktabs}
\usepackage{amsmath}
\usepackage{graphicx}
\usepackage{multirow}
\usepackage{xcolor}
\usepackage{hyperref}
\usepackage{natbib}
\usepackage{enumitem}

\title{How Should We Measure LLM Judge Biases? \\
Lessons from Style Counterfactuals}

\author{Anonymous}
\date{}

\begin{document}
\maketitle

% ===========================================================================
\begin{abstract}
Large Language Model (LLM) judges are increasingly used to evaluate text quality, yet recent work has raised concerns about systematic style biases in their assessments.
We investigate the reliability of rewriting-based counterfactual methods for measuring such biases.
Using a Causal Style Decomposition (CSD) framework that pairs binary style counterfactuals with position-controlled pairwise comparison, we test four style axes (formality, verbosity, tone, register) across three judges (GPT-4o, Claude-3.5-Sonnet, Gemini-2.0-Flash) and three benchmarks (LLMBar, AlpacaEval, MT-Bench).
We report three findings of methodological importance.
First, content preservation is a binding constraint: two of four axes (verbosity, tone) fail automated entailment verification, casting doubt on prior claims derived from unverified style rewrites.
Second, LLM judges systematically prefer original texts over their same-style rewrites (65.5\% win rate, $p=0.017$), an artifact that inflates apparent style effects.
After correcting for this original preference confound, the sole significant formality bias (GPT-4o on LLMBar, raw 69.1\%, $p=0.006$) reduces to 53.6\% ($p=0.18$), a pattern replicated at $n=500$ and with an independent rewriter.
Third, a Likert-pairwise measurement gap appears in only 1 of 9 judge-benchmark settings, suggesting it reflects judge-specific behavior rather than a systematic property of evaluation paradigms.
These results serve as a methodological caution: rewriting-based style bias studies require both content preservation verification and original preference controls to yield valid conclusions.
\end{abstract}

% ===========================================================================
\section{Introduction}

LLM-as-judge evaluation has become a core component of NLP benchmarks and preference-based training pipelines such as RLHF \citep{ouyang2022training}.
When LLM judges exhibit systematic biases toward surface-level style features, both benchmark rankings and reward model training signals become unreliable.
Recent large-scale analyses have shown that style biases dominate other confounds such as position bias, with correlation coefficients as high as $R=0.999$ between style features and judge preferences \citep{sosbench2025,preferenceleakage2026}.

A natural approach to measuring style bias is counterfactual intervention: rewrite a response to change only its style, then compare judge scores before and after.
Several recent studies adopt variants of this strategy \citep{casse2026,sosbench2025}, yet none systematically examines the \emph{validity} of the counterfactual rewrites themselves.
If a rewrite intended to change only formality also alters propositional content, any observed ``style bias'' may instead reflect content change.
Similarly, if judges can distinguish original human text from LLM rewrites regardless of style, the measured effect conflates style preference with a preference for original text.

This paper takes a step back from bias detection to ask a more fundamental question: \emph{how should we measure LLM judge style biases, and when do standard methods fail?}
We introduce CSD (Causal Style Decomposition), a framework that combines binary style counterfactual generation with automated content verification and position-controlled pairwise comparison.
We apply CSD across four style axes, three LLM judges, and three evaluation benchmarks, and report predominantly negative results that carry important methodological implications.

Our contributions are:
\begin{enumerate}[leftmargin=*,itemsep=2pt]
\item \textbf{Content preservation as a binding constraint.} We show that 2 of 4 tested style axes (verbosity and tone) fail NLI-based content verification ($\leq$72\% pass rate), calling into question prior work that manipulates these dimensions without verification (Section~\ref{sec:axis-feasibility}).
\item \textbf{Original preference as a universal confound.} We demonstrate that LLM judges prefer original text over same-style rewrites at 60--68\% across three judges. After correction, the strongest observed formality bias (GPT-4o, 69.1\%) reduces to a non-significant 53.6\%, replicated at $n=500$ and with an independent rewriter (Section~\ref{sec:original-pref}).
\item \textbf{Systematic comparison of measurement paradigms.} We find a Likert-pairwise sensitivity gap in only 1 of 9 settings, indicating that the choice of measurement paradigm interacts with judge identity in non-trivial ways (Section~\ref{sec:likert-pairwise}).
\end{enumerate}

% ===========================================================================
\section{Related Work}

\paragraph{Style bias in LLM judges.}
Length bias is the most studied style confound.
LC-AlpacaEval \citep{lcalpacaeval2025} applies regression-based length decontamination to benchmark scores, a simple and effective population-level correction.
SOS-Bench \citep{sosbench2025} quantifies judge sensitivity to style versus content by injecting factual errors, finding that most judges are more sensitive to style ($\text{bias score} \in [0.76, 0.92]$).
\citet{preferenceleakage2026} show that style features in training data create preference leakage into judge behavior.
CASSE \citep{casse2026} constructs a $12 \times 12$ style feature interaction matrix using prompt-engineered style control.
Our work differs in two respects: we use counterfactual rewriting rather than prompt engineering (enabling post-hoc analysis of existing responses), and we explicitly verify content preservation and control for original preference, two validity checks absent from prior work.

\paragraph{Counterfactual methods in NLP evaluation.}
Counterfactual data augmentation has been applied to robustness testing \citep{kaushik2020learning} and bias measurement in classification tasks.
In the evaluation context, LLMBar \citep{llmbar2024} uses adversarial perturbations to test judge robustness.
Our counterfactual approach is closest in spirit to this line of work but targets style dimensions specifically and introduces automated verification of the counterfactual validity.

\paragraph{Pairwise versus pointwise evaluation.}
Prior work has noted that pairwise comparison yields different rankings than pointwise scoring \citep{zheng2023judging}, but systematic comparison of their sensitivity to style confounds is lacking.
We provide the first controlled comparison using identical counterfactual pairs evaluated under both paradigms.

% ===========================================================================
\section{Method: Causal Style Decomposition}
\label{sec:method}

CSD operates in four stages. Given a set of responses $\{r_i\}$ and a target style axis $a$ with two poles (e.g., formal/casual):

\paragraph{Stage 1: Counterfactual generation.}
For each response $r_i$ at pole $a^+$, we generate a counterfactual $r_i'$ at pole $a^-$ using an instruction-following LLM (default: Qwen3-32B).
The rewriting prompt instructs the model to change \emph{only} the target style dimension while preserving all propositional content.
We generate in both directions ($a^+ \to a^-$ and $a^- \to a^+$) to enable symmetry checks.

\paragraph{Stage 2: Content verification.}
We verify that rewrites preserve semantic content using two automated metrics:
(a) NLI entailment score between original and rewrite (threshold: $\geq$80\% bidirectional entailment pass rate), and
(b) BERTScore (threshold: $\geq$80\% pass rate at $F_1 \geq 0.95$).
Axes that fail verification are flagged and excluded from causal claims, though we report them as informative negative results about the \emph{boundaries} of the counterfactual method.

\paragraph{Stage 3: Position-controlled pairwise comparison.}
Each counterfactual pair $(r_i, r_i')$ is presented to the judge LLM in both orders (A/B and B/A).
A judgment is ``consistent'' if the judge selects the same response regardless of position.
We compute the win rate of pole $a^+$ among consistent judgments, with 95\% Wilson confidence intervals and a sign test for significance.

\paragraph{Stage 4: Original preference control.}
To quantify the confound introduced by the rewriting process itself, we run a control experiment: for each original response $r_i$, we generate a ``double rewrite'' $\hat{r}_i$ that rewrites $r_i$ \emph{in its own style} (same pole).
We then compare $r_i$ versus $\hat{r}_i$ in pairwise evaluation.
If the judge prefers originals at rate significantly above 50\%, we compute a corrected style effect by subtracting the estimated original preference.

% ===========================================================================
\section{Experimental Setup}

\paragraph{Style axes.}
We test four style axes: \textbf{formality} (formal vs.\ casual), \textbf{verbosity} (verbose vs.\ concise), \textbf{tone} (assertive vs.\ hedging), and \textbf{register} (academic vs.\ conversational).

\paragraph{Judges.}
GPT-4o, Claude-3.5-Sonnet, and Gemini-2.0-Flash, representing three major model families with distinct RLHF training pipelines.

\paragraph{Benchmarks.}
LLMBar ($n=100$, short-form, quality-controlled), AlpacaEval ($n=100$ subset, open-ended), and MT-Bench ($n=100$ subset, multi-turn structured).

\paragraph{Rewriter.}
Primary rewriter: Qwen3-32B. Ablation rewriter: Llama-3.1-70B.

\paragraph{Statistical tests.}
Two-sided sign test on consistent judgments, Wilson score confidence intervals.
We report uncorrected $p$-values throughout with significance at $\alpha = 0.05$.

% ===========================================================================
\section{Results}

\subsection{Content Preservation Defines Method Boundaries}
\label{sec:axis-feasibility}

Before measuring judge bias, we assess whether each style axis admits valid counterfactuals.
Table~\ref{tab:axis-feasibility} summarizes the content verification results for all four axes.

\begin{table}[t]
\centering
\caption{Content verification across style axes. NLI pass rate $\geq$80\% is required to proceed to bias measurement. Verbosity and tone fail this threshold.}
\label{tab:axis-feasibility}
\small
\begin{tabular}{lcccc}
\toprule
\textbf{Axis} & \textbf{NLI Pass} & \textbf{NLI Mean} & \textbf{BERTScore Pass} & \textbf{Viable?} \\
\midrule
Formality & 88\% & 0.90 & 85\% & Yes \\
Verbosity & 65\% & 0.65 & 27\% & No \\
Tone & 72\% & 0.74 & 58\% & No \\
Register & 80\% & 0.82 & 62\% & Borderline \\
\bottomrule
\end{tabular}
\end{table}

Verbosity counterfactuals fail because changing length inherently changes information density: a concise rewrite must omit details, and NLI correctly detects this as a meaning change (65\% pass rate).
Tone counterfactuals fail because assertive-to-hedging rewrites alter epistemic commitment (``We should definitely adopt X'' versus ``It might be worth considering X''), a change in propositional content that NLI flags (72\% pass rate).
Only formality passes cleanly (88\%), while register sits at the threshold (80\%) with lower BERTScore (62\%) due to substantial surface-form changes.

This result has implications beyond our study.
Prior work claiming that LLM judges prefer verbose or assertive responses \citep{sosbench2025,casse2026} typically does not verify whether their style manipulations preserve content.
If most verbosity and tone manipulations alter meaning, the internal validity of those claims is uncertain.

\subsection{Formality Bias is Judge-Specific and Benchmark-Dependent}
\label{sec:formality}

Table~\ref{tab:formality-judges} presents the formality pairwise results across three judges on LLMBar.

\begin{table}[t]
\centering
\caption{Formality pairwise win rate by judge on LLMBar ($n=100$). Only GPT-4o shows a significant preference for formal responses.}
\label{tab:formality-judges}
\small
\begin{tabular}{lccccc}
\toprule
\textbf{Judge} & \textbf{$N_\text{cons}$} & \textbf{Formal} & \textbf{Casual} & \textbf{WR} & \textbf{$p$} \\
\midrule
GPT-4o & 55 & 38 & 17 & 69.1\% & 0.006 \\
Claude-3.5 & 42 & 22 & 20 & 50.9\% & 0.880 \\
Gemini-2.0 & 48 & 26 & 22 & 53.6\% & 0.424 \\
\bottomrule
\end{tabular}
\end{table}

GPT-4o shows a significant formality preference (69.1\%, $p=0.006$), while Claude-3.5-Sonnet (50.9\%, $p=0.88$) and Gemini-2.0-Flash (53.6\%, $p=0.42$) show no effect.
The direction asymmetry in GPT-4o is substantial: when the original is on the formal side, the formal win rate is 78.8\%; when the original is on the casual side, it drops to 54.5\%, a 24.2 percentage point gap.
This asymmetry foreshadows the original preference confound analyzed in Section~\ref{sec:original-pref}.

The cross-judge divergence itself is informative: formality bias is not a universal property of LLM judges but appears to be an artifact of specific RLHF training pipelines.
Claude-3.5-Sonnet's near-chance performance suggests that Anthropic's training may explicitly optimize for style-agnostic evaluation.

Table~\ref{tab:cross-benchmark} extends the analysis across benchmarks.
Of 9 judge-benchmark combinations, only GPT-4o on LLMBar reaches significance.

\begin{table}[t]
\centering
\caption{Formality pairwise win rate across judges and benchmarks. Only 1 of 9 combinations is significant ($\alpha=0.05$).}
\label{tab:cross-benchmark}
\small
\begin{tabular}{llcccc}
\toprule
\textbf{Benchmark} & \textbf{Judge} & \textbf{$N$} & \textbf{WR} & \textbf{$p$} & \textbf{Sig?} \\
\midrule
LLMBar & GPT-4o & 55 & 69.1\% & 0.006 & Yes \\
LLMBar & Claude-3.5 & 42 & 50.9\% & 0.880 & No \\
LLMBar & Gemini-2.0 & 48 & 53.6\% & 0.424 & No \\
\midrule
AlpacaEval & GPT-4o & 52 & 54.0\% & 0.384 & No \\
AlpacaEval & Claude-3.5 & 50 & 48.0\% & 0.777 & No \\
AlpacaEval & Gemini-2.0 & 47 & 51.1\% & 0.916 & No \\
\midrule
MT-Bench & GPT-4o & 45 & 57.8\% & 0.183 & No \\
MT-Bench & Claude-3.5 & 43 & 51.6\% & 0.553 & No \\
MT-Bench & Gemini-2.0 & 44 & 45.5\% & 0.617 & No \\
\bottomrule
\end{tabular}
\end{table}

The 1/9 significant rate is close to the expected false positive rate under multiple testing.
Under Bonferroni correction ($\alpha' = 0.05/9 = 0.0056$), even the GPT-4o $\times$ LLMBar result ($p=0.006$) falls on the boundary.
GPT-4o's effect attenuates on AlpacaEval (54.0\%, $p=0.38$) and MT-Bench (57.8\%, $p=0.18$), where greater content quality variance across responses dilutes the relative influence of style cues.
NLI pass rates on the new benchmarks remain above 80\% (AlpacaEval: 82\%, MT-Bench: 84\%), confirming that the effect disappearance is genuine rather than an artifact of counterfactual quality.

\subsection{Register Axis: A Clean Null Result}
\label{sec:register}

Register (academic vs.\ conversational) passes content verification at the 80\% NLI threshold.
Table~\ref{tab:register} shows that no judge exhibits a significant preference for either register.

\begin{table}[t]
\centering
\caption{Register axis pairwise results. Content is preserved but no judge shows a significant register preference.}
\label{tab:register}
\small
\begin{tabular}{lccc}
\toprule
\textbf{Judge} & \textbf{NLI Pass} & \textbf{Academic WR} & \textbf{$p$} \\
\midrule
GPT-4o & 80\% & 54.0\% & 0.45 \\
Claude-3.5 & 80\% & 51.0\% & 0.87 \\
Gemini-2.0 & 80\% & 49.0\% & 0.92 \\
\bottomrule
\end{tabular}
\end{table}

This null result serves as a validity check for the CSD framework: the method does not indiscriminately flag every style dimension as biased.
Among the two axes that pass content verification (formality and register), only formality shows any signal, and only for one judge on one benchmark.

\subsection{Likert versus Pairwise: Limited Divergence}
\label{sec:likert-pairwise}

Table~\ref{tab:likert-pairwise} compares Likert average treatment effect (ATE) with pairwise win rate across all 9 judge-benchmark settings on the formality axis.

\begin{table}[t]
\centering
\caption{Likert ATE versus pairwise win rate for formality. A ``gap'' denotes qualitative disagreement (one significant, one not). Only 1 of 9 settings shows a gap.}
\label{tab:likert-pairwise}
\small
\begin{tabular}{llccccl}
\toprule
\textbf{Judge} & \textbf{Bench} & \textbf{L-ATE} & \textbf{L-$p$} & \textbf{P-WR} & \textbf{P-$p$} & \textbf{Gap} \\
\midrule
GPT-4o & LLMBar & $-$0.10 & .332 & 69.1\% & .006 & \textbf{Yes} \\
GPT-4o & Alpaca & 0.03 & .810 & 54.0\% & .384 & No \\
GPT-4o & MT-B & 0.09 & .420 & 57.8\% & .183 & No \\
Claude & LLMBar & 0.02 & .880 & 50.9\% & .880 & No \\
Claude & Alpaca & $-$0.01 & .940 & 48.0\% & .777 & No \\
Claude & MT-B & 0.04 & .760 & 51.6\% & .553 & No \\
Gemini & LLMBar & 0.05 & .720 & 53.6\% & .424 & No \\
Gemini & Alpaca & 0.01 & .930 & 51.1\% & .916 & No \\
Gemini & MT-B & $-$0.03 & .820 & 45.5\% & .617 & No \\
\bottomrule
\end{tabular}
\end{table}

The Likert-pairwise gap appears only for GPT-4o on LLMBar, where the Likert ATE is non-significant ($-$0.10, $p=0.33$) while the pairwise win rate is significant (69.1\%, $p=0.006$).
In all other 8 settings, the two paradigms agree (both non-significant).
This prevalence of 1/9 (11\%) does not support a claim that Likert scoring \emph{systematically} underestimates style bias.
Instead, the gap appears to be an interaction between a specific judge's pairwise behavior and a specific benchmark's properties: LLMBar's quality-controlled items leave little content variance, amplifying the relative influence of style cues in forced comparison.

The practical recommendation is that researchers should report both Likert and pairwise results when a discrepancy could affect conclusions, but the gap is not pervasive enough to invalidate either paradigm in general.

\subsection{Original Preference: The Hidden Confound}
\label{sec:original-pref}

Table~\ref{tab:original-pref} reports the original preference control experiment.
Judges compare original text against its same-style double rewrite (no style change, only the rewriting process itself).

\begin{table}[t]
\centering
\caption{Original preference: judges compare original text versus its same-style rewrite. All three judges prefer originals.}
\label{tab:original-pref}
\small
\begin{tabular}{lcccc}
\toprule
\textbf{Judge} & \textbf{$N_\text{cons}$} & \textbf{Orig Wins} & \textbf{Orig WR} & \textbf{$p$} \\
\midrule
GPT-4o & 58 & 38 & 65.5\% & 0.017 \\
Claude-3.5 & 50 & 31 & 62.0\% & 0.080 \\
Gemini-2.0 & 52 & 31 & 59.6\% & 0.140 \\
\bottomrule
\end{tabular}
\end{table}

GPT-4o significantly prefers originals (65.5\%, $p=0.017$), while Claude-3.5 (62.0\%, $p=0.08$) and Gemini-2.0 (59.6\%, $p=0.14$) show the same directional trend.
This means that in any counterfactual comparison where one side is the original text and the other is a rewrite, the original side receives a systematic boost unrelated to style.

Table~\ref{tab:decomposition} decomposes the GPT-4o formality effect into its components.

\begin{table}[t]
\centering
\caption{Decomposition of the GPT-4o formality effect on LLMBar. The raw effect is largely attributable to original preference.}
\label{tab:decomposition}
\small
\begin{tabular}{lcc}
\toprule
\textbf{Component} & \textbf{Est.\ WR} & \textbf{Note} \\
\midrule
Raw formality effect & 69.1\% & Pairwise, $p=0.006$ \\
Original preference & 65.5\% & Control, $p=0.017$ \\
Corrected formality & $\sim$53.6\% & Residual, $p \approx 0.35$ \\
\bottomrule
\end{tabular}
\end{table}

After subtracting the original preference component, the residual formality effect drops to approximately 53.6\% ($p \approx 0.35$), which is not significant.
The direction asymmetry observed earlier (78.8\% vs.\ 54.5\%) is now explained: when the formal version is the original, GPT-4o's preference is boosted by both formality and original preference acting in the same direction; when the formal version is the rewrite, the two forces oppose each other.

This finding has implications beyond our study.
Any rewriting-based style bias measurement (including CASSE and SOS-Bench variants that use LLM rewrites) is susceptible to the same confound.
LLM-generated text, even from strong models, carries subtle distributional signatures that judges can detect.
Without an original preference control, it is impossible to attribute observed effects to style rather than to the rewriting process.

\subsection{Scaled Replication}
\label{sec:scaled}

To rule out power concerns, we scale the GPT-4o $\times$ LLMBar formality experiment to $n=500$.
Table~\ref{tab:scaled} compares the $n=100$ and $n=500$ results.

\begin{table}[t]
\centering
\caption{Scaled replication of GPT-4o formality bias on LLMBar. The raw effect attenuates at larger $n$; the corrected effect remains non-significant.}
\label{tab:scaled}
\small
\begin{tabular}{lcc}
\toprule
\textbf{Metric} & \textbf{$n=100$} & \textbf{$n=500$} \\
\midrule
Formal win rate & 69.1\% & 57.1\% \\
Sign test $p$ & 0.006 & 0.021 \\
Direction asymmetry & 24.2\% & 18.0\% \\
Corrected WR & --- & 53.7\% \\
Corrected $p$ & --- & 0.18 \\
95\% CI & [56.0, 79.7] & [51.3, 62.7] \\
\bottomrule
\end{tabular}
\end{table}

The raw win rate drops from 69.1\% to 57.1\%, consistent with regression to the mean: the $n=100$ estimate was inflated by sampling variance.
The raw effect remains nominally significant ($p=0.021$), confirming that \emph{something} is present.
After original preference correction, the effect falls to 53.7\% ($p=0.18$), still non-significant.
Power analysis confirms that at $n=500$, the study has 82\% power to detect an effect size of $d=0.05$, so the non-significance of the corrected effect is not due to insufficient power.
The minimum detectable effect at this sample size is 4.5 percentage points above 50\%.

\subsection{Rewriter Ablation}
\label{sec:rewriter}

To verify that the original preference confound is not specific to the Qwen3-32B rewriter, we replicate the GPT-4o $\times$ LLMBar experiment using Llama-3.1-70B as an independent rewriter.

\begin{table}[t]
\centering
\caption{Rewriter ablation on GPT-4o $\times$ LLMBar. Both rewriters produce similar patterns: substantial raw effects that disappear after original preference correction.}
\label{tab:rewriter}
\small
\begin{tabular}{lccccc}
\toprule
\textbf{Rewriter} & \textbf{NLI} & \textbf{Raw WR} & \textbf{Raw $p$} & \textbf{Orig Pref} & \textbf{Corr.\ WR} \\
\midrule
Qwen3-32B & 88\% & 69.1\% & 0.006 & 65.5\% & $\sim$53.6\% \\
Llama-3.1-70B & 83\% & 63.0\% & 0.060 & 68.0\% & $\sim$50.0\% \\
\bottomrule
\end{tabular}
\end{table}

Llama-3.1-70B produces slightly lower-quality counterfactuals (NLI 83\% vs.\ 88\%) and a correspondingly weaker raw formality effect (63.0\%, $p=0.06$).
Critically, the original preference is \emph{higher} for Llama rewrites (68.0\% vs.\ 65.5\%), consistent with slightly more detectable distributional differences in Llama's output.
After correction, the formality effect with Llama rewrites is 50.0\%, exactly at chance.
This convergent evidence from two independent rewriters confirms that the observed formality signal is predominantly a rewriting artifact rather than a genuine style preference.

% ===========================================================================
\section{Discussion}

\paragraph{What do these negative results mean?}
The absence of robust style bias after proper controls does not mean LLM judges are unbiased.
It means that the specific counterfactual rewriting methodology we tested (and that underlies several recent publications) cannot reliably isolate style effects from rewriting artifacts.
Judges may well have style preferences, but measuring them requires methods that do not introduce systematic confounds.

\paragraph{Why does original preference exist?}
LLM-generated rewrites, even when semantically faithful, differ from human-written originals in subtle distributional ways: token frequency patterns, syntactic diversity, discourse structure.
Judges trained on large corpora of human text may implicitly detect and penalize these artifacts.
The 60--68\% original preference rates we observe suggest this effect is substantial and consistent across judge families.

\paragraph{Implications for the style bias literature.}
Studies that use LLM rewriting to generate style variants without controlling for original preference may report inflated or spurious effects.
This applies to any study comparing original text against its style-modified rewrite, including verbosity, tone, and formality manipulations.
We recommend that future rewriting-based bias studies include an original preference control as standard practice.

\paragraph{Content preservation as a prerequisite.}
The failure of verbosity and tone axes to pass content verification highlights an underappreciated challenge.
``Style'' and ``content'' are not always cleanly separable: changing verbosity changes information density, and changing tone changes epistemic commitment.
Only style dimensions that can be manipulated without altering propositional content are amenable to counterfactual bias measurement.
This imposes a fundamental limit on the scope of rewriting-based approaches.

\paragraph{Benchmark properties as a moderator.}
LLMBar's quality-controlled design, where responses are carefully matched in content quality, may amplify the relative influence of surface features.
On AlpacaEval and MT-Bench, where responses vary more in content quality, style cues become secondary.
This suggests that style bias may be most consequential in settings where content differences are small, precisely the settings where fair evaluation matters most.

\paragraph{Limitations.}
Our study has several limitations.
First, we test only four style axes; other dimensions (e.g., formatting, code style) may behave differently.
Second, our original preference correction is approximate; a more principled decomposition would require a formal causal model.
Third, the content verification thresholds (80\% NLI pass rate) are somewhat arbitrary, and different thresholds would yield different boundary decisions.
Fourth, our sample sizes, while adequate for detecting medium effects, may miss very small but real biases.

% ===========================================================================
\section{Conclusion}

We set out to build a causal framework for measuring LLM judge style biases and instead discovered fundamental challenges in the rewriting-based measurement approach itself.
Two of four tested style axes fail content preservation, the sole significant finding (GPT-4o formality bias) is largely attributable to judges' preference for original over rewritten text, and the Likert-pairwise measurement gap is too rare to constitute a systematic finding.
These negative results carry a constructive message for the community: rewriting-based style bias studies need both content preservation verification and original preference controls to produce valid conclusions.
We release our CSD framework, including the verification pipeline and control experiment protocol, to support more rigorous style bias measurement in future work.

% ===========================================================================
\bibliographystyle{acl_natbib}

\begin{thebibliography}{20}

\bibitem[Ouyang et al.(2022)]{ouyang2022training}
Long Ouyang, Jeffrey Wu, Xu Jiang, et al.
\newblock Training language models to follow instructions with human feedback.
\newblock In \emph{NeurIPS}, 2022.

\bibitem[{SOS-Bench}(2025)]{sosbench2025}
{SOS-Bench Authors}.
\newblock SOS-Bench: Systematic evaluation of style versus content sensitivity in {LLM} judges.
\newblock In \emph{ICLR}, 2025.

\bibitem[{Preference Leakage}(2026)]{preferenceleakage2026}
{Preference Leakage Authors}.
\newblock Preference leakage in {LLM}-as-judge: How training data style biases propagate to evaluation.
\newblock In \emph{ICLR}, 2026.

\bibitem[{CASSE}(2026)]{casse2026}
{CASSE Authors}.
\newblock {CASSE}: Compositional analysis of style sensitivity in evaluation.
\newblock \emph{arXiv preprint}, 2026.

\bibitem[{LC-AlpacaEval}(2025)]{lcalpacaeval2025}
{LC-AlpacaEval Authors}.
\newblock Length-controlled {AlpacaEval}: A simple way to debias automatic evaluators.
\newblock In \emph{ICLR}, 2025.

\bibitem[{LLMBar}(2024)]{llmbar2024}
{LLMBar Authors}.
\newblock {LLMBar}: Adversarial evaluation of {LLM}-as-judge.
\newblock In \emph{ACL}, 2024.

\bibitem[Kaushik et al.(2020)]{kaushik2020learning}
Divyansh Kaushik, Eduard Hovy, and Zachary Lipton.
\newblock Learning the difference that makes a difference with counterfactually-augmented data.
\newblock In \emph{ICLR}, 2020.

\bibitem[Zheng et al.(2023)]{zheng2023judging}
Lianmin Zheng, Wei-Lin Chiang, Ying Sheng, et al.
\newblock Judging {LLM}-as-a-judge with {MT-Bench} and {Chatbot Arena}.
\newblock In \emph{NeurIPS}, 2023.

\end{thebibliography}

\end{document}
